% Page 2: abstract, contents and the seven claims (the reader's entry point).
\pagenumbering{arabic}
\setcounter{page}{2}

\begin{center}
  \begin{minipage}{0.92\linewidth}
  \small
  \textbf{Abstract.}
  Alert thresholds on non-stationary signals are often allowed to learn from recent
  data. This report asks \emph{what} they should learn from. A threshold that admits
  only the windows it already judged normal can only fall, and the falling threshold
  floods the operator with alerts; admitting every window instead absorbs sustained
  faults into the new normal. We treat recalibration as an explicit decision-time
  action, chosen alongside \textsc{alert}, \textsc{hold} and \textsc{defer}, whose
  training data is selected independently of the threshold's own decisions and kept
  within a band anchored to the calibration threshold. Eleven policies are replayed on
  one common score trace per stream, and alert workload, decision coverage, delay and
  event recall are reported separately. On two SKAB valve streams and seven synthetic
  stress cases we prove the failure, measure its cost, and find the new controller no
  worse than a fixed threshold on SKAB but not yet shown to be better on real data. We
  state the experiment that would confirm or refute the direction.
  \end{minipage}
\end{center}

% Compact two-column contents: section level only, no paragraph gaps.
{\hypersetup{linkcolor=ink}\setcounter{tocdepth}{1}\setlength{\parskip}{0pt}%
 \makeatletter\renewcommand\tableofcontents{\@starttoc{toc}}%
 \renewcommand*\l@section[2]{\@dottedtocline{1}{0pt}{1.6em}{\small #1}{\small #2}\vspace{1pt}}\makeatother
 \noindent\textbf{Contents}\par\vspace{2pt}
 \begin{multicols}{2}\tableofcontents\end{multicols}}

\glancebox{%
  \textbf{Our research direction and its seven claims}\par\smallskip
  \emph{Recalibrating an alert threshold should be a guarded decision-time action whose
  training data is chosen neither by that same threshold nor unconditionally, and every
  alternative should be judged on one common score trace.}
  Tags say how each claim is supported: \Sproved{} by proof, \Smeasured{} on saved
  results, \Simplemented{} in tested code, \Sproposed{} not yet tested.
  \begin{enumerate}[label=\textbf{C\arabic*}, leftmargin=2.2em, itemsep=2pt, topsep=4pt]
    \item \Sproved{} We prove that a threshold which learns only from windows it judged
      normal can never rise (Proposition~\ref{prop:ratchet}).
    \item \Smeasured{} We measure the cost on two SKAB streams: the same single event is
      detected as with the fixed threshold, at $16.2\times$ (valve1) and $3.9\times$
      (valve2) the alert windows.
    \item \Smeasured{} We measure that the scorer, not the policy, limits detection:
      only 6/100 and 13/98 in-event windows exceed the calibration threshold.
    \item \Simplemented{} We implement recalibration as one of four actions, bounded to
      $[\theta_0, 4\theta_0]$, and compare it with seven baselines and three ablations on
      the identical trace.
    \item \Smeasured{} On SKAB the controller is identical to the fixed threshold on both
      streams: a no-harm check, not evidence of benefit.
    \item \Smeasured{} Under synthetic stress it cuts false-alert time on a benign step
      (16 vs 597) and still catches later faults, but fails on a benign ramp and on a
      fault shaped like a step, as declared in advance.
    \item \Sproposed{} We propose, and do not claim, a benefit on real benign regime
      changes; a predeclared experiment with three kill criteria decides it
      (Section~\ref{sec:conclusion}).
  \end{enumerate}
  \textbf{Closest prior work.} Each ingredient is already published: the failure in
  FITNESS~\citep{sankararaman2022fitness}, abstention in Sun et
  al.~\citep{sun2024online} and Perini \& Davis~\citep{perini2023rejection},
  workload-based evaluation in Hu~\citep{hu2026workload}. What we did not find in the
  nine closest papers is their combination on one fixed score trace
  (Table~\ref{tab:occupation}). The strongest case against us is in
  Section~\ref{sec:limitations}.
}
